\clearpage % Rozdziały zaczynamy od nowej strony.
\section{Ewaluacja, strojenie progów i analiza efektywności (FinOps)}
\label{chap:ewaluacja-i-finops}

% W niniejszym rozdziale przedstawiono wyniki eksperymentów badawczych przeprowadzonych na systemie \texttt{s3-dedup-engine}. Skupiono się na procesie optymalizacji i strojenia progów decyzyjnych modelu percepcji, analizie wydajności potoku przetwarzania oraz kwantyfikacji oszczędności finansowych w oparciu o model FinOps.

% -------------------------------------------------------------------------
\subsection{Metodologia eksperymentu i przygotowanie danych}
\label{sec:metodologia-eksperymentu}

% TODO: Opisać zbiory danych wykorzystane do testów (np. pobrane za pomocą scripts/fetch_dataset.py) oraz strukturę eksperymentu badawczego.
% \begin{itemize}
%     \item Charakterystyka testowego zbioru danych graficznych (zróżnicowanie pod kątem rozdzielczości, stopnia kompresji JPEG oraz transformacji geometrycznych).
%     \item Kryteria podziału na zbiory referencyjne i zapytania testowe (\textit{query vs. database}).
% \end{itemize}

% -------------------------------------------------------------------------
\subsection{Strojenie progów podobieństwa (Threshold Tuning)}
\label{sec:threshold-tuning}

% TODO: Opisać wyniki eksperymentów ze strojenia progu cosinusowego w oparciu o notatnik scripts/tune_threshold.ipynb.
% \begin{itemize}
%     \item Analiza wpływów wartości granicznych miary podobieństwa cosinusowego (\textit{Cosine Similarity}) na skuteczność detekcji duplikatów.
%     \item Krzywe precyzji i pełności (\textit{Precision-Recall}) oraz dobór optymalnego progu operacyjnego minimalizującego liczbę fałszywych alarmów (\textit{False Positives}) przy jednoczesnym zachowaniu wysokiej wykrywalności kopii.
% \end{itemize}

% -------------------------------------------------------------------------
\subsection{Wydajność potoku deduplikacyjnego (Pipeline Performance)}
\label{sec:wydajonsc-potoku}

% TODO: Przedstawić wyniki pomiarów narzutu czasowego (latency) i obciążenia zasobów.
% \begin{itemize}
%     \item Analiza skuteczności strategii wczesnego wyjścia (\textit{Early Exit}): jaki odsetek plików jest eliminowany na etapie wstępnym (nazwa, pierwszy chunk, pełne SHA-256) bez konieczności angażowania modelu AI.
%     \item Porównanie czasu operacji I/O w systemie hybrydowym (Go + Python via gRPC) względem tradycyjnego, pełnego przetwarzania sekwencyjnego.
% \end{itemize}

% -------------------------------------------------------------------------
\subsection{Analiza finansowa i optymalizacja kosztów chmurowych (FinOps)}
\label{sec:analiza-finops}

% TODO: Opisać symulację oszczędności kosztowych w ekosystemie Amazon S3.
% \begin{itemize}
%     \item Kwantyfikacja redukcji zajętości przestrzeni dyskowej dzięki eliminacji duplikatów binarnych (1:1) oraz percepcyjnych (\textit{upgrade} jakościowy poprzez \textit{InfoGain}).
%     \item Symulacja oszczędności na operacjach API (wywołania typu \texttt{PUT}, \textit{GET}, \textit{LIST}) oraz transferze danych wychodzących w oparciu o cennik Amazon S3.
% \end{itemize}